\section{Relation to realization and automata theory}
The polynomial matrix upper construction realizes an entire legal
terminal density matrix, not merely one acceptance probability. In
particular, enlarging its decoder set to the full probability simplex
after tomography changes the problem. The extreme-point proof requires
the legal image of $\mathcal D_d$.

Brodsky--Pippenger \cite[Theorems~3.3 and 3.6; Lemma~3.5]{BP} characterize
bounded-error measure-once quantum language recognition by group
languages, discuss a free-group word problem, and transfer cut-point
recognition to probabilistic automata. Those statements compare languages
and cut-points; our exact profile constrains the full terminal output and
every cut of a horizon-specific stochastic machine. Ambainis--Freivalds
\cite{AF57} likewise provide an important automata-complexity comparison,
not a substitute for that legal-output quantifier.

The geometric ingredients are classical smooth homogeneous-space and
polytope arguments; compare \cite{Hall,Szarek57}. The least-rank theorem
makes the cap estimate uniform over all conditional-centroid directions,
including singular orbit types. The quaternion argument gives complete
finite-input freeness and stabilizer proofs. Bourgain--Gamburd
\cite[Theorem~1]{BG2012} is a separate deep input for the actual displayed
alphabet's noisy converse. Neither that spectral theorem nor classical
free subgroup constructions are counted as new realization theorems.

Positive realization and invariant-cone methods \cite{BF04,BR} explain
why a common positive lift, rather than unrelated one-cut factorizations,
is necessary. Hidden-process realization \cite{Heller58} concerns
consistent output processes; the present model is a terminal controlled
interface, and makes no assertion about adaptive observations. These
comparisons specify different mathematical questions, not an exhaustive
priority determination.

There remain two quantitative issues even for the displayed rational
example. The gap between $N/\log N$ and $N$ at fixed positive error has
not been closed, and the transition between the proved exponentially
small-error regime and a fixed-error regime has not been classified.
The constructed rational upper bound is effective without a spectral-gap
constant. No practical total-complexity claim is made for general
compact-group minimization.
